\documentclass[11pt]{article}

\usepackage[letterpaper,margin=0.88in]{geometry}
\usepackage{amsmath,amssymb,amsfonts}
\usepackage{array,booktabs,longtable,tabularx}
\usepackage{enumitem}
\usepackage{titlesec}
\usepackage{fancyhdr}
\usepackage{lastpage}
\usepackage{hyperref}
\usepackage{xcolor}
\usepackage{graphicx}

\setlength{\headheight}{14pt}
\hypersetup{
  colorlinks=true,
  linkcolor=black,
  urlcolor=black,
  citecolor=black
}

\pagestyle{fancy}
\fancyhf{}
\lhead{White Casimir Physical Readouts}
\rhead{Apparatus Path Paper}
\cfoot{\thepage\ of \pageref{LastPage}}

\titleformat{\section}{\large\bfseries\uppercase}{\thesection.}{0.5em}{}
\titleformat{\subsection}{\normalsize\bfseries}{\thesubsection}{0.5em}{}
\setlist[itemize]{nosep,leftmargin=1.2em}
\setlist[enumerate]{nosep,leftmargin=1.4em}
\renewcommand{\arraystretch}{1.16}

\newcommand{\um}{\,\mu\mathrm{m}}
\newcommand{\sbr}{\mathrm{SBR}}
\newcommand{\fp}{\mathrm{FP}}
\newcommand{\Tc}{T_{\mathrm{c}}}
\newcommand{\df}{\delta f/f}
\newcommand{\pathlocal}[1]{\texttt{#1}}

\title{\bfseries From White-Style Casimir Shells to Physical Readout Paths:\\
High-Q, Superconducting, and Differential-Pressure Apparatus Options}
\author{
Daniel S. Goldman\\
Independent researcher\\
\href{https://orcid.org/0000-0003-2835-3521}{ORCID: 0000-0003-2835-3521}
}
\date{\today}

\begin{document}
\maketitle

\begin{abstract}
White et al.'s sphere-in-cylinder Casimir calculation created a useful split in
the experimental question. One part is geometric: whether a custom Casimir
boundary condition can organize a shell-like negative vacuum-stress proxy near
the region visually compared with an Alcubierre source shell. A recent
role-resolved audit found that this shell-morphology layer survives several
proxy and perturbation tests, while the direct metric/timing readout layer does
not survive calibrated scale and electromagnetic-background accounting. The
result is not a warp-readout program. It is a narrower physical program:
identify whether a shell-coded Casimir boundary response can be transduced by
ordinary precision apparatus strongly enough to be measured and separated from
material, patch-potential, thermal, vibration, and readout backgrounds.

This paper starts from that narrowed claim boundary and translates the Stage 5
readout ladder into apparatus paths. The three remaining physical branches are
a high-Q electromagnetic cavity frequency/phase shift, a superconducting
resonator frequency/phase shift, and a paired differential-pressure cell. The
first two are the strongest surviving branches under clean-control priors. The
differential-pressure cell is more direct as force metrology but survives only
under unusually strong patch-potential and common-mode control. A material or
impedance readout is retained as a calibration/control branch, not as evidence
for a shell response. The old central timing branch closes by amplitude.

The proposed continuation is therefore a Casimir-boundary metrology program.
A positive result would not show practical warp behavior or a freely usable
negative-energy fluid. It would show that a finite shell-coded boundary
condition produces a repeatable renormalized-stress or boundary-mode response
that survives physical controls in a modern apparatus. That would still be
valuable: it would make the White-style geometry a precision test bed for
finite-geometry Casimir physics, surface-patch systematics, high-Q resonant
transduction, superconducting vacuum-boundary effects, and possible Casimir
device engineering.
\end{abstract}

\section{Claim Boundary}

White, Vera, Han, Bruccoleri, and MacArthur used worldline numerics on custom
Casimir geometries and reported that one sphere-in-cylinder configuration
generated a negative energy-density morphology qualitatively similar to a
two-dimensional Alcubierre shell \cite{White2021}. The compact geometry is a
$1\um$ diameter sphere centered in a nominal $4\um$ diameter cylinder. The
important fact for this paper is not the visual resemblance alone. The useful
fact is that the geometry separates three questions:
\begin{enumerate}
\item Does the boundary condition organize a shell-like Casimir response?
\item Is the calibrated stress-energy scale remotely competitive with a metric
source requirement?
\item Is there a laboratory observable that can read the shell response through
ordinary electromagnetic and material backgrounds?
\end{enumerate}

The prior role-resolved audit answered these in different directions. The
source-placement layer was positive in the limited proxy sense: shell
morphology remained concentrated in the annular sphere-wall region and was
controlled by the sphere and cylinder geometry. The metric-source and central
timing layers were negative: the calibrated scale was many orders below an
Alcubierre-type source demand, and the central readout path did not carry a
recoverable physical signal. The synthetic discrimination layer added one more
condition. If a shell-coded observable already reaches near-background
amplitude, a geometry schedule can separate it from nuisance templates at about
\(\sbr=1.2\). The schedule is a discriminator, not an amplitude generator.

The Stage 5 run therefore changed the research question. The question is no
longer whether the White geometry gives a direct warp or metric readout. In
this harness, it does not. The question is whether the shell-coded Casimir
boundary response can be coupled into a practical physical transducer strongly
enough to pass an apparatus gate. The word ``physical'' is doing real work
here: the branch must specify a sensor, a noise floor, dominant nuisance
channels, calibration controls, a geometry modulation schedule, and a failure
criterion.

\section{Negative Energy Is Not Just a Label}

The first conceptual question is whether these readouts would say anything
about negative energy in quantum field theory, or whether ``negative energy''
is only bookkeeping. The careful answer is that QFT negative energy is
operational but constrained. In renormalized stress-energy accounting, the
local expectation value of energy density can be lower than the reference
vacuum in some region. The Casimir effect is one familiar expression of this:
boundaries alter the mode structure and produce measurable force, pressure,
and energy changes. Ford's review emphasizes both sides of the point: local
negative energy densities are allowed in QFT, while quantum inequalities and
compensating positive energy restrict their magnitude, duration, and use
\cite{Ford2009}.

So a successful White-style physical readout would not merely attach a name to
a calculation. If it survived geometry modulation, material substitution,
surface-potential mapping, dummy-cell subtraction, and readout-nuisance
controls, it would provide operational evidence that the finite boundary
condition changed a measurable physical quantity in the direction predicted by
renormalized vacuum-stress or Casimir-boundary theory. Depending on the branch,
that quantity could be a force/pressure, a mechanical frequency shift, a
cavity eigenfrequency shift, a microwave phase shift, or a superconducting
resonator response.

But the same result would not elevate the system into an exotic-matter or
warp-source device. The Stage 3 scale comparison already placed the
gravitational/source magnitude far beneath representative metric demand. A
positive high-Q or superconducting readout would mean: the shell-coded
boundary condition is physically readable as a Casimir-boundary response. It
would not mean: the negative energy has become an unconstrained material,
extractable fuel, or engineered stress tensor suitable for macroscopic metric
control. The strongest interpretation remains QFT-metrological, not
propulsion-metrological.

\section{Stage 5 Inputs}

The Stage 5 priority physical run used the Stage 3 shell proxy and Stage 4
synthetic discrimination threshold as input. It evaluated five candidate
readout families across physical backend cases. The run produced 33 candidate
scenario gates, 27 physical-backend rows, 237600 synthetic observations, and
19800 fits. The declared gate threshold inherited from Stage 4 was
\(\sbr=1.2\), with ordinary-EM-only false-positive control near
\(\fp<0.05\).

\begin{table}[h]
\centering
\caption{Stage 5 gate counts for the priority physical closeout run.}
\label{tab:gate_counts}
\small
\begin{tabular}{lr}
\toprule
Gate status & Count \\
\midrule
Candidate survives for experiment & 11 \\
Closed physical scale & 14 \\
Closed nuisance degeneracy & 2 \\
Closed amplitude & 3 \\
Background calibration only & 3 \\
\bottomrule
\end{tabular}
\end{table}

\begin{table}[h]
\centering
\caption{Highest-ranked Stage 5 survivors and near-survivors.}
\label{tab:stage5_survivors}
\small
\begin{tabularx}{\linewidth}{>{\raggedright\arraybackslash}p{0.30\linewidth}lrr>{\raggedright\arraybackslash}X}
\toprule
Branch & Scenario & Pred. \(\sbr\) & \(\fp\) & Dominant nuisance \\
\midrule
High-Q cavity shift & optimistic clean & 19.34 & 0.040 & loss channel \\
Superconducting resonator & optimistic clean & 17.76 & 0.0367 & surface loss \\
Superconducting resonator & optimistic balanced & 12.10 & 0.040 & surface loss \\
High-Q cavity shift & optimistic matched & 11.88 & 0.0267 & loss channel \\
Superconducting resonator & optimistic lossy & 7.94 & 0.0433 & surface loss \\
High-Q cavity shift & optimistic lossy & 6.07 & 0.0433 & loss channel \\
Superconducting resonator & nominal clean & 5.46 & 0.0267 & surface loss \\
Differential pressure cell & optimistic quiet & 5.31 & 0.0267 & patch potential \\
Superconducting resonator & nominal balanced & 3.59 & 0.0333 & surface loss \\
High-Q cavity shift & nominal clean & 2.22 & 0.0133 & loss channel \\
Differential pressure cell & optimistic balanced & 2.02 & 0.030 & patch potential \\
\bottomrule
\end{tabularx}
\end{table}

The central timing reference is retained in the ladder as a closure check. Its
optimistic physical row has predicted \(\sbr\sim4.1\times10^{-59}\), and the
three central timing rows all close by amplitude. The material impedance branch
has predicted \(\sbr=0\) by construction because it is a background and
calibration channel, not a shell measurement. The surviving physical program is
therefore concentrated in resonant boundary readouts, with differential
pressure as a harder but conceptually direct force channel.

\section{Experimental Architecture}

The viable apparatus cannot be a single static measurement. It has to be a
geometry-modulated differential measurement, because the nuisance templates
share many controls with the shell template. The experimental object is a
finite boundary cell with controlled shell-gap modulation, sphere/cylinder
geometry, wall thickness, readout coupling, and matched dummy geometry. A
single absolute shift is not enough. The evidence is the vector of shifts
across a modulation schedule.

The common architecture has five parts. First, a shell cell implements the
sphere-cylinder or a manufacturable surrogate geometry with a known gap,
surface roughness, conductivity, coating thickness, and alignment tolerance.
Second, at least one dummy cell carries the same material, wiring, thermal
mass, and readout chain while suppressing the shell overlap. Third, a geometry
schedule changes the coordinates that the Stage 3/Stage 4 templates identified
as discriminating: gap, cylinder radius, sphere size or surrogate radius,
offset, wall thickness, length, and readout overlap. Fourth, an independent
surface and material metrology layer maps roughness, residual potential,
finite conductivity, coating thickness, loss tangent, and thermal expansion.
Fifth, the readout is fitted against both a shell template and ordinary
background templates, with an EM-only false-positive test.

This architecture is close to established Casimir metrology practice. Decca,
Aksyuk, and Lopez review MEMS/NEMS force measurements as a natural platform
for Casimir sensing \cite{Decca2011}. Garcia-Sanchez et al. used metallized
high-Q silicon nitride nanomembranes and in situ surface-potential measurement
to control residual electrostatic forces \cite{GarciaSanchez2012,
GarciaSanchez2013}. Bimonte et al. measured differential Casimir forces from
0.2 to \(8\um\) while explicitly treating surface roughness, edge effects, and
patch potentials \cite{Bimonte2021}. These are not optional niceties for the
White-style follow-on; they are the minimum machinery required to keep a
Casimir claim from becoming a material-background claim.

\section{Path A: High-Q Cavity Shift}

The high-Q cavity branch treats the shell cell as a perturbation to an
electromagnetic resonator. The observable is a fractional frequency or phase
shift,
\begin{equation}
  y_{\mathrm{cav}} = \left(\frac{\delta f}{f}\right)_{\mathrm{meas}},
\end{equation}
measured over a geometry schedule. In the Stage 5 physical backend, the best
high-Q case used a shell fractional frequency shift of
\(9.40\times10^{-11}\) against a background fractional rms of
\(4.86\times10^{-12}\), giving predicted \(\sbr=19.34\). The strongest
nominal clean case gave \(\sbr=2.22\). Matched and lossy cases rapidly lose
margin, which is exactly the point: the branch is viable only when the mode
overlap and loss model are strong enough.

The physical device would use a cavity mode with high electric or magnetic
field overlap in the shell-sensitive region and a companion null or low-overlap
mode. A vector network analyzer, Pound-style lock, or phase-locked loop tracks
the resonant frequency and quality factor. The fit model has the form
\begin{equation}
  y_{\mathrm{cav}}(q) =
  a_{\mathrm{shell}}T_{\mathrm{shell}}(q)
  + a_{\mathrm{cond}}T_{\mathrm{cond}}(q)
  + a_{\mathrm{therm}}T_{\mathrm{therm}}(q)
  + a_{\mathrm{wg}}T_{\mathrm{wg}}(q)
  + a_{\mathrm{mix}}T_{\mathrm{mix}}(q)
  + a_{\mathrm{loss}}T_{\mathrm{loss}}(q)
  + \epsilon ,
\end{equation}
where \(q\) denotes the geometry setting and the \(T_i\) are shell and
nuisance templates. The cavity branch does not win by measuring a larger
force directly. It wins by using resonance to make a small boundary-coded
perturbation readable.

The required machinery is therefore mode-engineering machinery: a validated
finite-element eigenmode solver; a measured mode-overlap factor; independent
thermal expansion tracking; dummy cavities; controlled wall and conductor
loss; and a way to reverse, detune, or suppress the shell-overlap mode while
leaving the material and readout chain intact. Without these controls, a
frequency shift is easy to obtain and hard to interpret. With them, a positive
result would imply that the finite Casimir boundary condition couples to an
electromagnetic eigenmode in a reproducible, geometry-coded way.

\section{Path B: Superconducting Resonator Shift}

The superconducting resonator branch is the most interesting physically
because it sits at the intersection of Casimir boundary conditions and
cryogenic electrodynamics. The observable is again a fractional frequency or
phase shift, but now the transducer is a superconducting resonator, microwave
optomechanical structure, or superconducting boundary device. Stage 5 gave the
best superconducting clean case \(\sbr=17.76\), with nominal clean at
\(\sbr=5.46\) and nominal balanced at \(\sbr=3.59\). The conservative and
lossy rows show the main risk: surface loss, vortex loss, quasiparticle loss,
kinetic inductance, and material-control channels can remain too close to the
shell readout.

This branch has external motivation beyond the local harness. Perez et al.
designed a NEMS system to probe Casimir energy corrections to superconducting
condensation energy by measuring superconducting transition changes in a
controlled cavity \cite{Perez2020}. de Jong et al. report measurement of the
Casimir force between superconducting objects through the nonlinear dynamics
of a superconducting drum resonator in a microwave optomechanical system
\cite{DeJong2025}. Gurevich reviews the loss physics that must be confronted
in superconducting resonators, including quasiparticles, dielectric loss,
residual surface resistance, trapped vortices, and limits on high-Q operation
\cite{Gurevich2023}. These sources make the branch concrete: this is the kind
of machine that can plausibly read a tiny boundary perturbation, and also the
kind of machine whose internal loss channels can counterfeit one.

The clean experiment would use two resonator classes. The first is a material
and loss reference resonator with the same film, substrate, wiring, field
environment, and thermal cycle but no shell-sensitive Casimir geometry. The
second is the shell-coupled resonator. The schedule should include temperature
sweeps across and below \(\Tc\), controlled magnetic-field history for vortex
loading, drive-power sweeps for quasiparticle and TLS response, and geometry
modulation for shell template recovery. A strong positive result would have to
remain locked to the shell geometry while changing differently from kinetic
inductance, surface resistance, vortex, quasiparticle, and material-control
templates.

If it worked, the physical implication would be narrower and sharper than a
warp claim. It would indicate that superconducting boundary conditions or
cryogenic resonant structures can transduce finite Casimir geometry changes
into a reproducible microwave observable. Practical utility could include
surface-loss diagnostics, Casimir-calibrated resonator sensors, vacuum-gap
metrology, and tests of how superconducting electrodynamics enters finite
geometry Casimir forces.

\section{Path C: Differential Pressure Cell}

The differential-pressure branch is conceptually the most direct. Instead of
asking a resonator to magnify a boundary perturbation, it tries to measure a
force or pressure contrast between matched cells. For ideal parallel plates,
the pressure scale is
\begin{equation}
  P_{\mathrm{C}}(a) = -\frac{\pi^2\hbar c}{240a^4},
\end{equation}
which is about \(1.3\times10^{-3}\,\mathrm{Pa}\) at \(a=1\um\), before real
geometry, material, finite-temperature, and roughness corrections. Stage 5's
optimistic pressure rows sit in this same experimental neighborhood: the
optimistic quiet row used a shell pressure of \(7.18\times10^{-4}\,\mathrm{Pa}\)
and passed with \(\sbr=5.31\). The nominal quiet row fell to
\(\sbr=0.35\), and patchy or conservative rows close quickly.

The pressure branch therefore requires a paired-cell design rather than a
single membrane. The two cells should share temperature, vibration, readout
electronics, electrode materials, and membrane mechanics while moving through
opposite shell-gap modulations. Patch-potential mapping is central, not
auxiliary. Bimonte et al.'s long-range force measurement is directly relevant
because it treats experimental uncertainty, patch potentials, roughness, edge
effects, and theory comparison without fitting parameters \cite{Bimonte2021}.
Stange et al.'s commercial MEMS platform shows how a practical force sensor
can be converted into a Casimir metrology platform with piconewton-scale
resolution and electrostatic calibration \cite{Stange2019}. The
nanomembrane/Kelvin-probe work shows why surface-potential control must be
built into the apparatus \cite{GarciaSanchez2012}.

The branch's positive implication would be the cleanest statement of force
physics: a shell-coded boundary geometry produces a pressure contrast that
tracks the template after common-mode subtraction. Its risk is also the
cleanest: patch potentials and ordinary electrostatic forces are naturally
large compared with the desired residual. In the Stage 5 ladder, pressure is
not dismissed, but it is the least forgiving survivor.

\section{Controls That Are Evidence-Bearing}

The material impedance branch should be kept, but only as a control. It is a
first-class competitor, not a shell readout. The most dangerous mistake would
be to interpret any geometry-correlated impedance, capacitance, or loss change
as evidence for a Casimir shell. Material response is exactly what must be
measured so that it can be regressed out or used to close the claim.

\begin{table}[h]
\centering
\caption{Control requirements by branch.}
\label{tab:controls}
\small
\begin{tabularx}{\linewidth}{>{\raggedright\arraybackslash}p{0.24\linewidth}>{\raggedright\arraybackslash}p{0.29\linewidth}>{\raggedright\arraybackslash}X}
\toprule
Branch & Required controls & Fail condition \\
\midrule
High-Q cavity & Mode-overlap solver, dummy cavity, thermal monitor, conductor-loss model, waveguide/mode-mixing checks & Frequency shift follows loss, thermal, or waveguide template more strongly than shell template \\
Superconducting resonator & Material-only resonator, vortex history, quasiparticle and drive-power sweeps, surface-loss model, \(\Tc\) sweep & Shift remains degenerate with kinetic inductance, surface loss, vortex, quasiparticle, or material-control channels \\
Differential pressure & Matched dummy cell, opposite gap modulation, Kelvin-probe map, roughness prior, vibration and thermal common-mode rejection & Residual pressure follows patch/electrostatic or membrane-drift template \\
Material impedance & Four-terminal calibration, capacitance/impedance mapping, dummy readout path & Treated as shell evidence rather than background calibration \\
Central timing & Matched central path and EM-only schedule & Closed by amplitude unless a new transduction model changes scale by many orders \\
\bottomrule
\end{tabularx}
\end{table}

The evidence-bearing form is a multi-template fit, not an isolated
instrumental shift. A plausible shell readout must pass four conditions:
\begin{enumerate}
\item It exceeds the Stage 4 recovery threshold under the same geometry
schedule, or a more stringent validated schedule.
\item EM-only and material-only synthetic or physical control datasets do not
generate false positives above the declared rate.
\item A physically measured nuisance template is not more correlated with the
data than the shell template.
\item The amplitude and sign remain stable under dummy-cell, material-swap,
and geometry-reversal tests.
\end{enumerate}

\section{Comparison with White et al.}

The proposed program keeps White et al.'s most durable contribution and drops
the fragile one. White et al. showed that a custom Casimir geometry can produce
a shell-like vacuum-response morphology in a worldline calculation
\cite{White2021}. The prior audit did not erase that. It made the statement
more precise: the shell morphology is a boundary-geometry result, while the
metric-source and central-timing readouts fail under calibrated scale and
background accounting.

The continuation proposed here changes the observable. White-style central or
transit language asks whether the apparatus creates a direct metric-like
effect readable along the center. The Stage 5 answer is no for the modeled
apparatus. The high-Q, superconducting, and pressure branches instead ask
whether the same shell-coded boundary geometry can be read through ordinary
Casimir metrology. That is a much more conservative question, but it is also
physically buildable.

The distinction matters for interpretation. A positive resonator or pressure
result would strengthen the claim that the finite shell geometry is a real
physical Casimir boundary condition, not a plotting artifact. It would not
resurrect the direct warp-readout interpretation. The best narrative is:
White et al. identified a provocative shell-forming Casimir geometry; the
audit separated morphology from metric source scale; Stage 5 identified the
physical readout branches worth testing as Casimir-boundary metrology.

\section{Build Plan}

The practical next program should be staged so that each step can close the
line cleanly.

\subsection{Phase 0: apparatus digital twin}

Build branch-specific simulations before fabricating the shell cell. For
the cavity branch, solve eigenmodes and perturbative frequency shifts for the
actual geometry, including conductor thickness, ports, readout feed, and wall
loss. For superconducting resonators, include kinetic inductance, surface
resistance, quasiparticle, vortex, TLS, and material-control terms. For
pressure cells, include membrane mechanics, patch maps, roughness, electrostatic
calibration, and edge corrections. The output of Phase 0 is a measured or
simulated template matrix that can be passed through the Stage 4/Stage 5
synthetic gates.

\subsection{Phase 1: known-geometry calibration}

Do not begin with the White shell cell. First reproduce a standard
sphere-plane, parallel-plate, trench, or membrane Casimir measurement at the
chosen scale. The purpose is not novelty. It is to prove the apparatus can
recover an ordinary Casimir signal while controlling residual electrostatics,
roughness, thermal drift, and readout artifacts. This phase should use the
Decca/Bimonte/Garcia-Sanchez style controls: electrostatic calibration,
Kelvin-probe surface-potential mapping where possible, separation metrology,
and explicit error budgets \cite{Decca2011,Bimonte2021,GarciaSanchez2012}.

\subsection{Phase 2: dummy and material controls}

Construct dummy geometries with the same materials and readout chain but
suppressed shell overlap. Run the full geometry schedule on shell and dummy
cells. This is where the material impedance branch earns its keep. A large
geometry-correlated material response is not a surprise; it is the main
background map.

\subsection{Phase 3: shell-cell measurement}

Run the shell cell under the full schedule. The acceptance gate should be
predeclared: predicted and measured \(\sbr\), false-positive rate, nuisance
correlations, mode/loss separation, and branch-specific amplitude bounds. If
the high-Q and superconducting branches fail after measured template updates,
the line should close as a physical readout program. If one branch survives,
the next paper should be an apparatus design and uncertainty-budget paper, not
a metric interpretation paper.

\section{Physical Implications and Utility}

If all three branches fail under measured apparatus templates, the physical
conclusion is still useful. It would say that the White-style shell morphology
is a real proxy-level organization of Casimir boundary response, but that
available transducers do not lift it through ordinary material and
electromagnetic backgrounds. That would close the current line at the
metrology gate.

If the pressure branch survives, the implication is a direct finite-geometry
force result. It would show that a shell-coded Casimir boundary condition can
produce a pressure contrast that remains after patch and common-mode controls.
Practical utility would be strongest in precision force metrology, MEMS/NEMS
calibration, and surface-patch diagnostics.

If the high-Q cavity branch survives, the implication is resonant
electromagnetic transduction. The shell boundary condition would be readable as
a mode shift rather than as a direct force. Practical utility could include
Casimir-assisted cavity sensors, mode-overlap diagnostics for microcavities,
and precision tests of finite-geometry vacuum-boundary calculations.

If the superconducting branch survives, the implication is the most
interesting for condensed-matter/QFT overlap. It would suggest that
superconducting boundary conditions, resonator loss physics, and Casimir
geometry can be separated well enough to measure a shell-coded vacuum-boundary
effect. Practical utility could include cryogenic gap metrology,
superconducting surface-loss diagnostics, and improved understanding of how
superconducting electrodynamics enters Casimir forces.

None of these outcomes would imply practical warp propulsion. The value is
smaller but real: a controlled way to ask how engineered boundaries shape
renormalized vacuum stress and how those changes can be read by precision
machines.

\section{Code, Data, and Source Availability}

The local source bundle for this paper is stored at
\begin{quote}
\small
\path{white_casimir_intersection/sources}
\end{quote}
and is documented in
\begin{quote}
\small
\path{white_casimir_intersection/sources/README.md}.
\end{quote}

The Stage 5 priority physical closeout used here is stored under:
\begin{quote}
\small
\path{/media/kir/9CDCBD3EDCBD140C/Research/}\\
\path{white_casimir_stage5_readout_ladder/stage5_priority_physical_closeout_20260530}
\end{quote}
The compact files used for the tables and claims are:
\begin{itemize}
\item \path{summary.json}
\item \path{gates/gate_summary.csv}
\item \path{candidates/readout_candidates.csv}
\item \path{backends/physical_backend_sweep.csv}
\end{itemize}

\section{Conclusion}

The current state of the White Casimir line is narrow but not empty. The
central timing and metric-readout interpretation is closed in the present
harness. The shell morphology is still useful as a finite Casimir boundary
configuration. Stage 5 says that the strongest physical continuations are
resonant readouts: a high-Q cavity and a superconducting resonator. A
differential-pressure cell remains a possible force readout only under strong
patch and common-mode control.

The next honest question is whether one of those physical branches can survive
measured apparatus templates. If the answer is yes, the result would be a
valuable Casimir-boundary readout, not a warp signal. If the answer is no, the
research line closes cleanly: White-style shell morphology exists at the proxy
level, but no physically robust readout remains after ordinary backgrounds and
controls.

\begin{thebibliography}{12}

\bibitem{White2021}
H. White, J. Vera, A. Han, A. R. Bruccoleri, and J. MacArthur,
``Worldline numerics applied to custom Casimir geometry generates
unanticipated intersection with Alcubierre warp metric,''
\emph{European Physical Journal C} \textbf{81}, 677 (2021).
\url{https://doi.org/10.1140/epjc/s10052-021-09484-z}.

\bibitem{Alcubierre1994}
M. Alcubierre,
``The warp drive: hyper-fast travel within general relativity,''
\emph{Classical and Quantum Gravity} \textbf{11}, L73--L77 (1994).
\url{https://doi.org/10.1088/0264-9381/11/5/001}.

\bibitem{Ford2009}
L. H. Ford,
``Negative Energy Densities in Quantum Field Theory,''
arXiv:0911.3597 (2009).
\url{https://arxiv.org/abs/0911.3597}.

\bibitem{Decca2011}
R. S. Decca, V. Aksyuk, and D. Lopez,
``Casimir force in micro and nano electro mechanical systems,''
in \emph{Casimir Physics}, Lecture Notes in Physics \textbf{834} (2011).
\url{https://tsapps.nist.gov/publication/get_pdf.cfm?pub_id=906575}.

\bibitem{GarciaSanchez2012}
D. Garcia-Sanchez, K. Y. Fong, H. Bhaskaran, S. Lamoreaux, and H. X. Tang,
``Casimir Force and In Situ Surface Potential Measurements on
Nanomembranes,''
\emph{Physical Review Letters} \textbf{109}, 027202 (2012).
\url{https://doi.org/10.1103/PhysRevLett.109.027202}.

\bibitem{GarciaSanchez2013}
D. Garcia-Sanchez, K. Y. Fong, H. Bhaskaran, S. Lamoreaux, and H. X. Tang,
``Casimir probe based upon metallized high Q SiN nanomembrane resonator,''
arXiv:1301.7025 (2013).
\url{https://arxiv.org/abs/1301.7025}.

\bibitem{Stange2019}
A. Stange, M. Imboden, J. Javor, L. K. Barrett, and D. J. Bishop,
``Building a Casimir metrology platform with a commercial MEMS sensor,''
\emph{Microsystems \& Nanoengineering} \textbf{5}, 14 (2019).
\url{https://doi.org/10.1038/s41378-019-0054-5}.

\bibitem{Perez2020}
D. J. Perez, A. Stange, R. Lally, L. Barrett, M. Imboden, A. Som,
D. Campbell, V. Aksyuk, and D. Bishop,
``A system for probing Casimir energy corrections to the condensation
energy,''
\emph{Microsystems \& Nanoengineering} \textbf{6}, 1 (2020).
\url{https://doi.org/10.1038/s41378-020-00221-2}.

\bibitem{Bimonte2021}
G. Bimonte, B. Spreng, P. A. Maia Neto, G.-L. Ingold,
G. L. Klimchitskaya, V. M. Mostepanenko, and R. S. Decca,
``Measurement of the Casimir Force between 0.2 and 8 um: Experimental
Procedures and Comparison with Theory,''
\emph{Universe} \textbf{7}, 93 (2021).
\url{https://doi.org/10.3390/universe7040093}.

\bibitem{Gurevich2023}
A. Gurevich,
``Tuning microwave losses in superconducting resonators,''
\emph{Superconductor Science and Technology} \textbf{36}, 073002 (2023).
\url{https://doi.org/10.1088/1361-6668/acc6e3}.

\bibitem{DeJong2025}
M. H. J. de Jong, E. Korkmazgil, L. Banniard, M. A. Sillanp\"a\"a,
and L. Mercier de L\'epinay,
``Measurement of the Casimir force between superconductors,''
arXiv:2501.13759 (2025).
\url{https://arxiv.org/abs/2501.13759}.

\end{thebibliography}

\end{document}
